\scanpage{90}

\mspage{\unsure{5|S|2}}

\[
G=\underline{\mathrm{Aut}}_S(F)
\]

(schéma en groupes affine sur $S$), alors $\operatorname{gr}(F)$ en tant que foncteur-fibre (en oubliant sa graduation) est défini par un $G$-torseur $P$, \inserted{par $\operatorname{gr}(F)\simeq F^P$, $P=\underline{\mathrm{Isom}}_S(F,\operatorname{gr}F)$}

De plus, la graduation sur \deleted{$\unsure{G|\otimes}$} le foncteur \deleted{sur} $\operatorname{gr}(F)$ est \unsure{défini|définie} par un hom

\[
G_{m_S}\xrightarrow{i}\deleted{\unsure{G'=}}G^P=G'.
\]

Toutes ces constructions sont can. et compatibles avec changement de base.

Parmi les isom de $F$ sur $\operatorname{gr}(F)$, on distingue ceux qui respectent la filtration \inserted{et induisent l'identité sur les gradués associés}. Ils forment (pour $S$ variable) un sous-foncteur $Q$ de $P$, qui est un \inserted{quasi-}torseur à g sous le sous-groupe $H'$ de $G'^{\deleted{P}}$ des autom qui respectent la filtration \inserted{canonique} de $F'=\operatorname{gr}F$\inserted{ et induisent l'identité sur $\operatorname{gr}\unsure{F'}$}, et un \inserted{quasi-}torseur à droite sous le sous-groupe analogue $H$ de $G$.

\margin{On suppose $\mathcal M$ de type fini}

Cependant, \inserted{pour voir que c'est un torseur,} il faut s'assurer que loc. \unsure{fppf|fpqc} $Q$ admet une section, donc qu'il existe un isom \uline{respectant les filtrations} de $F$ et $F'$ et induisant l'identité sur les gr associés.
% archive sheet number: 89
